\section{Marco teórico y trabajos relacionados}

Una cosa es reconocer la urgencia de procesar datos en tiempo real y otra muy distinta es diseñar la infraestructura capaz de soportarlo. Si en la sección anterior establecimos el "porqué" del cambio hacia arquitecturas orientadas a eventos (EDA), aquí es necesario diseccionar el "cómo". La literatura técnica reciente sugiere que la solución no es una "bala de plata", sino una convergencia de tres frentes: sacar el cómputo de la nube hacia el borde, dotar de significado semántico a los eventos y orquestar flujos de trabajo que van más allá de la simple mensajería.

\subsection{La descentralización obligatoria: Inteligencia en el Borde}
La premisa de enviar todos los datos a la nube para su análisis se ha vuelto insostenible, tanto por costos de ancho de banda como por latencia. La tendencia actual, defendida por autores como Baz et al. \cite{baz2019video} y Salamon et al. \cite{salamon2018gateway}, es transformar los dispositivos IoT y gateways en filtros inteligentes. La idea es simple pero potente: el video o la lectura del sensor no debe salir del dispositivo a menos que constituya un evento relevante.

Esta filosofía exige repensar incluso cómo capturamos la realidad. Un ejemplo fascinante es el trabajo de Coretti, Varile y Bertaina \cite{coretti2025neuromorphic}, quienes se alejan de las cámaras tradicionales que graban cuadros completos. Su enfoque con visión neuromórfica procesa únicamente los cambios de luminosidad en los píxeles (eventos físicos), reduciendo drásticamente la carga computacional necesaria para tareas críticas como evitar colisiones espaciales. En la tierra, vemos un paralelismo en los sistemas de tráfico: Perera et al. \cite{perera2024traffic} y Ciechomski et al. \cite{ciechomski2018soccer} demuestran que combinar detectores ligeros (como YOLO) con reglas deterministas locales es mucho más efectivo que intentar que una IA remota "entienda" todo el contexto de un partido de fútbol o una intersección vial.

\subsection{El problema de la "Torre de Babel" en los datos}
Ahora bien, mover datos rápido no sirve de nada si los sistemas no se entienden entre sí. Roest y Florkowski \cite{roest2017challenges} tocan una fibra sensible al hablar de la "brecha semántica": el abismo que existe entre un píxel cambiando de color y el concepto de "vehículo estacionado ilegalmente". Los sensores, por defecto, son tontos; no saben de contexto.

Para solucionar esto, investigadores como Phuttharak y Loke \cite{phuttharak2023smartcity} proponen capas de traducción semántica que estandarizan los datos heterogéneos antes de que entren al bus de eventos. Pero la estandarización es solo el primer paso. Bruns y Dunkel \cite{bruns2014fusion} llevan la discusión hacia la inferencia de estados. Su investigación sobre coordinación de ambulancias resalta que un sistema reactivo no puede ser amnésico; necesita mantener un contexto (stateful processing) mediante Máquinas de Estados Finitos (FSM) para deducir, a partir de eventos simples de GPS, si una unidad está "disponible", "en ruta" o "atendiendo", sin intervención humana.

\subsection{Orquestación: Más allá del Pub/Sub}
Finalmente, llegamos a la columna vertebral del sistema. Si bien herramientas como Apache Kafka son el estándar para mover estos volúmenes de datos, la forma de estructurar los servicios ha evolucionado. Kul y Tashiev \cite{kul2021microservices} ilustran la potencia de descomponer tareas monolíticas (analizar un video) en microservicios paralelos y pequeños (uno mira el color, otro la velocidad), que se comunican asíncronamente. Esto permite una velocidad de respuesta que el procesamiento secuencial jamás podría igualar.

Sin embargo, no todo es color de rosa. En sectores donde un error cuesta millones o implica sanciones legales, como el financiero, la "coreografía" libre de eventos puede volverse caótica e inauditable. Hentrich y Pachmann \cite{adekunle2024financial} sugieren un enfoque híbrido prudente: usar motores de procesos (como Camunda) para orquestar los flujos críticos sobre la capa de eventos, garantizando que siempre haya un responsable del estado de la transacción. Esta cautela resuena con la advertencia de Lazzari y Farias \cite{lazzari2023hidden}, quienes nos recuerdan que la mayor debilidad de estas arquitecturas distribuidas sigue siendo la dificultad para rastrear errores y monitorear la salud del sistema cuando las cosas salen mal.